\documentclass[11pt]{article}
\usepackage[margin=1in]{geometry}
\usepackage{amsmath,amssymb,amsthm,mathtools,booktabs,hyperref}
\setlength{\emergencystretch}{3em}
\newtheorem{theorem}{Theorem}
\newtheorem{lemma}[theorem]{Lemma}
\newtheorem{definition}[theorem]{Definition}
\newcommand{\SAT}{\operatorname{SAT}}
\newcommand{\gammats}{2\cos(\pi/7)}
\title{Main 6: Polynomial-Regime Self-Improvement and the WRAPPER Bridge}
\author{PNP research working manuscript}
\date{7 October 2026}
\begin{document}
\maketitle

\begin{abstract}
This manuscript isolates the closest published analogue of the recursive
WRAPPER, Williams's Circuit-SAT self-improvement theorem, and determines
exactly why it does not yet extend a bare assumption $P=NP$ to the joint
resource region needed here.  It then gives a source-checkable sufficient
theorem for a forward WRAPPER tree.  The recovered WRAPPER components prove
correctness, recursion, sharing and several accounting interfaces, but no
current construction proves the theorem's all-input quantitative premises.
The result is a sharper conditional route, not a proof that $P\ne NP$.
\end{abstract}

\section{Purpose and relation to the other spines}
\texttt{main4.tex} studies existence-level contraction of the joint
achievable region.  \texttt{main5.tex} studies explicit recursive SAT
interfaces.  This manuscript records the bridge suggested by the double
WRAPPER: an inner solver is represented inside a forward object, outer
variables are left exposed, and an outer solver or evaluator completes the
computation.  This is precisely the architecture of the ``OR trick'' in
Williams's self-improvement proof \cite{williams2024}.

Three statuses must remain separate.
\begin{description}
\item[] \textbf{Known theorem.} Williams proves self-improvement in the
large-circuit and parameterized regimes.
\item[] \textbf{Proved conditional bridge.} A forward source-queryable tree
with the joint bounds below would yield a forbidden SAT algorithm.
\item[] \textbf{Open construction.} No recovered general-input WRAPPER has
the required node, generation, depth, leaf-length and workspace bounds in one
realization.
\end{description}

\section{The published self-improvement theorem}
Let an $n$-input circuit have size $2^{\alpha n+o(n)}$.  Williams's
Theorem~1.1 states that, for fixed positive constants $\alpha\le\beta$, a
$2^{\beta n+o(n)}$-time Circuit-SAT algorithm on that large-circuit regime
implies a $2^{(\beta-\alpha)n+o(n)}$-time algorithm on
$2^{o(n)}$-size circuits \cite{williams2024}.  The proof repeatedly splits
variables, represents the current inner SAT algorithm by a circuit, and asks
Circuit-SAT on the resulting outer circuit.  It is a rigorous recursive plus
forward-wrapper construction.

Suppose only that $P=NP$ and that one fixed Circuit-SAT implementation takes
time $N^{K+o(1)}$.  On size $N=2^{\alpha n+o(n)}$ circuits this gives
$\beta=K\alpha$, so Theorem~1.1 yields
\[
 2^{(K-1)\alpha n+o(n)}
\tag{1}
\]
time on subexponential-size circuits.  Choosing a different fixed $\alpha$
for each desired exponential rate recovers arbitrarily small exponential
rates, which is consistent with $P=NP$.  It does not give a fixed polynomial
bound in the encoded length of every subexponential-size input.

\begin{lemma}[Fixed-$\alpha$ polynomial-regime gate]
Fix $K>1$ and $\alpha>0$.  For a family whose encoded circuit length is
$M=2^{r(n)n}$ with $r(n)\to0$, the bound in (1), expressed as a power of $M$,
has exponent
\[
 \frac{(K-1)\alpha}{r(n)}+o(1/r(n)).
\]
Consequently the published fixed-$\alpha$ theorem gives no finite uniform
polynomial exponent in $M$ on the entire $2^{o(n)}$ regime.
\end{lemma}
\begin{proof}
Substitute $n=\log_2(M)/r(n)$ into (1).  Since $r(n)\to0$, the leading
coefficient diverges.  Choosing $\alpha$ as a function of the particular
input length would require a new uniform theorem controlling its construction,
induction depth and error terms.  The cited theorem fixes $\alpha$ as a
constant before the asymptotics.
\end{proof}

The parameterized result makes the second obstruction equally explicit.
Williams's Lemma~4.1 composes size bounds
\[
 O(2^{ak}+n^b)\quad\hbox{and}\quad O(2^{ck}+n^d)
\]
into one with variable exponent $a'=adc/(ad+c)$ but size exponent
$b'=bd$.  Self-composition of a generic $n^K$ size term therefore produces
$n^{K^2}$ while improving the variable term.  Theorem~1.7 avoids this growth
because its premise supplies a bound $O(c^k+n^{1+\varepsilon})$ for one fixed
$c>1$ and for every $\varepsilon>0$.

\begin{lemma}[Hybrid-premise gate]
The two bounds supplied without a new WRAPPER,
\[
 n^{K+o(1)}\quad\hbox{and}\quad 2^k n^{1+o(1)},
\]
do not by themselves satisfy
\[
 \exists c>1\ \forall\varepsilon>0:\quad
 O(c^k+n^{1+\varepsilon}).
\tag{2}
\]
More precisely, for every fixed $c>1$ and $K>1$ there are arbitrarily small
$\varepsilon>0$ and parameter ratios $k=2\varepsilon\log_2n+O(1)$ for which
the exponent certified by the hybrid minimum is $1+2\varepsilon$, whereas
the two terms on the right of (2) have exponent at most $1+\varepsilon$.
\end{lemma}
\begin{proof}
Write $d=\log_2c$.  Choose $\varepsilon>0$ so small that
$1+2\varepsilon<K$ and $2d\varepsilon<1$.  At
$k=2\varepsilon\log_2n+O(1)$, the brute-force term is
$n^{1+2\varepsilon+o(1)}$, and it beats $n^K$.  Meanwhile
$c^k=n^{2d\varepsilon+o(1)}$ and the other term is
$n^{1+\varepsilon}$.  Their sum has exponent $1+\varepsilon$, which is
strictly smaller than the only exponent certified by the stated hybrid.
This is a failure of those two upper bounds to prove (2), not a lower bound
against another algorithm.
\end{proof}

Williams explicitly asks whether self-improvement can reach the $P=NP$
regime and says the present method has bottlenecks below subexponential time
\cite{williams2024}.  Thus the remembered elimination of $K$ is valid inside
the project's conditional shrinker recurrences, and Williams eliminates a
variable exponent under stronger size-sensitive premises, but neither result
derives the missing premise from bare $P=NP$.

\section{A complete conditional forward-tree bridge}
The following theorem is the cleanest current target for the nonrecursive
inner WRAPPER plus recursive outer WRAPPER.

\begin{definition}[Source-queryable forward tree]
On an explicit CNF $F$ of encoded length $N$, a deterministic constructor
defines a finite rooted AND/OR tree $\mathcal T_F$.  Each leaf is a standard
explicit CNF.  A depth-first evaluator can regenerate the next child and the
combiner state from the read-only source and its current stack.  Let
\begin{align*}
 Q(N)&=\text{worst-case number of tree nodes actually evaluated},\\
 D(N)&=\text{maximum evaluated root-to-leaf depth},\\
 G(N)&=\text{total construction, regeneration and combination time},\\
 R(N)&=\text{peak generator and combiner workspace},\\
 M(N)&=\text{maximum encoded length of an evaluated leaf}.
\end{align*}
The tree is exact when its Boolean value equals $\SAT(F)$ for every input.
All five functions refer to the same implementation and include unsuccessful
routes needed on worst-case SAT and UNSAT inputs.
\end{definition}

\begin{theorem}[FORWARD-TREE-BRIDGE1]
Assume $P=NP$.  Suppose one exact source-queryable forward-tree WRAPPER has
fixed constants $q,g$ and satisfies
\[
 Q(N)\le N^{q+o(1)},\qquad G(N)\le N^{g+o(1)},
\]
\[
 D(N)=N^{o(1)},\qquad R(N)=N^{o(1)},\qquad M(N)=N^{o(1)},
\]
with a depth-first frame of $N^{o(1)}$ bits.  Then SAT has one deterministic
algorithm with
\[
 T(N)=N^{\max\{q,g\}+o(1)},\qquad S(N)=N^{o(1)}.
\]
If $\max\{q,g\}<\gammats$, this contradicts the model-matched Williams
time-space product theorem \cite{williams2010}.
\end{theorem}
\begin{proof}
Under $P=NP$, fix one exact SAT decider $A$ with time $m^{K+o(1)}$ and space
$m^{L+o(1)}$ for finite constants $K,L$.  Evaluate $\mathcal T_F$ depth
first, calling $A$ at a leaf.  Since $M(N)=N^{o(1)}$, every leaf call takes
$N^{o(1)}$ time and space.  The total leaf cost is therefore
$Q(N)N^{o(1)}=N^{q+o(1)}$.  Adding the fully charged $G(N)$ gives the time
bound.  At most one leaf and one root-to-current-node stack are live.  The
leaf, generator state and $D(N)$ frames together occupy $N^{o(1)}$ bits.
The endpoint excludes the resulting exponent pair when the strict inequality
holds.
\end{proof}

If generation cost $N^{h+o(1)}$ is only a per-node bound, the theorem must use
$q+h$ in place of $g$.  A short descriptor is not a small ordinary SAT
instance.  A generator for a virtually large leaf must charge every access,
and applying ordinary SAT to its descriptor would instead invoke a succinct
problem whose complexity can be much higher \cite{py1986}.

The tree condition is also substantive.  Formula/tree value has low-depth
evaluation structure, whereas evaluation of a general shared Boolean DAG is
the Circuit Value Problem and is complete for deterministic polynomial time
under logspace reductions \cite{ladner1975,buss1987}.  Sharing is useful only
with a proved evaluator bound for the resulting DAG class, or after charging
its unfolded tree size.  This is an interface guard, not a claim that every
shared WRAPPER requires large space.

There is a concrete intermediate target for the recursive version.  A
bounded-fan-in circuit of polynomial size and treewidth $w$ can be balanced
to formula depth $O(w^2\log N)$ \cite{jansensarma2010}.  Query-DAG results
likewise turn separator number $s$ into $2^{O(s\log N)}$ deterministic work
and $O(s\log N)$ oracle queries \cite{gharibianrudolph2023}.  Thus a WRAPPER
may retain sharing if it emits a source-checkable decomposition and the
resulting evaluator cost is charged.  Constant width controls dependency and
space, but the published generic polynomial-time transformations do not fix
an exponent below $\gammats$.  The remaining theorem must therefore bound
the actual evaluator time, not merely assert bounded treewidth or separator
number.  This alternative route separates the recursive-topology problem
from the still-open leaf-size and acquisition problem.

The selector literature marks another exact boundary.  A merge that must
return one of two inputs and preserve their disjunction is a p-selector;
p-selectivity together with self-reducibility characterizes polynomial-time
decidability \cite{buhrmantorenvliet1996}.  Under $P=NP$ such a selector
exists, but its polynomial exponent is uncontrolled.  A cheap selector would
solve the semantic choice problem, while its construction cost is precisely
part of the resource theorem still needed here.

\section{Where the aggregate $2^k/2^x$ calculation fits}
If a declared state population begins at $D_0=2^k$ and exact reductions have
factors $G_i$, then
\[
 D_r=\frac{2^k}{\prod_iG_i}=2^{k-x},\qquad
 x=\sum_i\log_2G_i.
\]
The condition $D_r\le N^d$ is exactly
$x\ge k-d\log_2N$.  For a subpolynomial residual population, it becomes
$x\ge k-o(\log N)$.  In FORWARD-TREE-BRIDGE1 this calculation may certify a
bound on $Q$ or $M$ only after a source-checkable map connects the counted
states to actual tree nodes or explicit leaf bits.  Counting nominal
assignments, hashes, prime labels, or merged route names does not establish
that connection.  This is the precise place where the original divide-by-
$2^x$ argument belongs.

For fixed-power leaves $M=N^{b+o(1)}$ instead, a hypothetical $N^K$ solver
contributes $q+bK$ to the time exponent.  The finite-depth theorem in
\texttt{main5.tex} removes the unknown contribution only after a repeatable
shrinking return is proved and the depth is chosen so that $b^dK$ fits the
endpoint margin.  FORWARD-TREE-BRIDGE1 is the limiting $b=o(1)$ form.

\section{What the latest constructions establish}
The double WRAPPER already supplies exact forward AND/OR semantics, standard
CNF leaves below its declared cutoff, strict recursive termination and
normalized-formula reuse.  The source-backed variants avoid retained residual
CNFs.  FRONTIER-JOINT1 supplies an exact future-boundary congruence, and
BATCHFRONTIER1 shares a suffix scan across blocks.  STREAMCOMP1 and
ORBITBATCH1 supply conditional charged ledgers for implicit access and shared
setup.  These are solved components, not wholly open placeholders.

The latest batched frontier implementation has, for sparse bounded-width
inputs with block exponent $\beta$ and table exponent $\omega$, the envelope
\[
 C=\max\{2-\beta,1+\omega\},\qquad
 \sigma=\max\{\beta,\omega\}.
\]
For every $\beta,\omega\ge0$, $C+\sigma\ge2$: if
$\beta\ge\omega$, the first branch $2-\beta$ plus $\beta$ gives two;
if $\omega>\beta$, the branch $2-\beta$ plus $\omega$ exceeds two.  The implementation's
chain control realizes the corresponding $N^{2-o(1)}$ product.  This retires
that full-suffix-scan access plan as an endpoint solution, not frontier
merging or the easy chain family.

A combined source-replay prototype now places propagation, frontier DP,
component AND, Shannon OR and collision-checked bounded sharing in one
implementation.  Its finite correctness tests are useful, but pigeonhole
controls exhibit rapid growth in recursion and source rereads.  Its Python
workspace instrumentation now charges simultaneously live component,
normalization, frontier, cache and ancestor descriptors.  The resulting
finite figures are accounting checks only: the dense assignment mask alone
is $\Theta(n)$ bits, and no asymptotic premise of
FORWARD-TREE-BRIDGE1 is inferred from those measurements.

Random bitstreams and prime lineage can audit identity or probabilistically
fingerprint replay.  The endpoint here is deterministic.  A favorable fixed
stream is advice unless it is generated uniformly; input bits can be reread;
and exact prime products carry their full bit length.  No recovered bitstream
construction supplies the five all-input bounds above.

The exact fingerprint audit in \texttt{main7.tex} sharpens this statement.
If cache hits are replay-verified, the digest length changes candidate work
but not the number of future-congruence classes, hence it does not bound the
tree-node function $Q(N)$.  If equal digests are merged without replay, the
claim that every digest fiber is one congruence class is already the missing
semantic contraction theorem.  Fixed four-, six- and eight-hex-digit outputs
are constants and cannot change a polynomial exponent.  A logarithmically
growing randomized fingerprint can reduce expected collision work, but a
fixed public map retains hostile collision fibers and does not meet this
deterministic all-input bridge.

\section{Reproducible arithmetic audit}
The companion program \texttt{polynomial\_regime\_bridge\_test.py} performs
635 exact rational checks in under one second, with a 25-second internal
deadline.  It checks the substitution $\beta=K\alpha$, 49 explicit witnesses
to the hybrid-premise gate, divergence of the fixed-$\alpha$ encoded-size
exponent, the BATCHFRONTIER1 minimum product of two on a $1/16$ grid, and the
disappearance of $K$ from the forward-tree ledger when the leaf exponent is
zero.  Raw output is in
\path{polynomial_regime_bridge_results.json}.  The analytic lemmas above,
not the finite grid, establish the symbolic claims.

\section{Current proof spine and exact open statement}
The strongest assembled route is
\[
\begin{aligned}
P=NP
&\Rightarrow \text{one fixed SAT decider with finite }(K,L)\\
&\overset{\textsc{Open}}{\Longrightarrow}
  \text{FORWARD-TREE-BRIDGE1 premises, or an ARC successor}\\
&\overset{\textsc{Proved/Exact conditional composition}}{\Longrightarrow}
  (C,\sigma)\text{ with }C+\sigma<\gammats\\
&\overset{\textsc{Known theorem}}{\Longrightarrow}\bot.
\end{aligned}
\]
Equivalently, one may replace the tree premises by a shared-DAG theorem with
a model-matched evaluator, or by the achievable-region finite-return premise
of \texttt{main4.tex}.  What remains open is one SAT-specific quantitative
implication, not correctness of the existing wrappers or the downstream
recurrence algebra:
\begin{center}
\fbox{\parbox{0.86\linewidth}{
Construct, on every explicit SAT input, one exact source-matched realization
whose total acquisition, branching, recursive leaf length, evaluation and
peak workspace jointly enter the Williams forbidden region; or prove the
corresponding existence implication between achievable resource pairs.
}}
\end{center}
Until this statement is proved, the three manuscripts give conditional proof
spines and falsifiable construction targets.  P versus NP remains open.

\begin{thebibliography}{9}
\bibitem{williams2024}
R. Ryan Williams,
\emph{Self-Improvement for Circuit-Analysis Problems},
STOC 2024, 1374--1385.
\url{https://people.csail.mit.edu/rrw/stoc-self-improvement.pdf}

\bibitem{williams2010}
R. Ryan Williams,
\emph{Alternation-Trading Proofs, Linear Programming, and Lower Bounds},
STACS 2010, Theorem 3.9.
\url{https://drops.dagstuhl.de/entities/document/10.4230/LIPIcs.STACS.2010.2494}

\bibitem{ladner1975}
R. E. Ladner,
\emph{The Circuit Value Problem is Log Space Complete for P},
SIGACT News 7(1), 1975.
\url{https://doi.org/10.1145/990518.990519}

\bibitem{buss1987}
S. R. Buss,
\emph{The Boolean Formula Value Problem is in ALOGTIME},
STOC 1987.
\url{https://mathweb.ucsd.edu/~sbuss/ResearchWeb/Boolean/paper.pdf}

\bibitem{py1986}
C. H. Papadimitriou and M. Yannakakis,
\emph{A Note on Succinct Representations of Graphs},
Information and Control 71(3), 1986.
\url{https://doi.org/10.1016/S0019-9958(86)80009-2}

\bibitem{jansensarma2010}
M. Jansen and J. Sarma,
\emph{Balancing Bounded Treewidth Circuits},
Theoretical Computer Science 2010.
\url{https://arxiv.org/abs/0910.1427}

\bibitem{gharibianrudolph2023}
S. Gharibian and D. Rudolph,
\emph{On Polynomially Many Queries to NP or QMA Oracles},
Theory of Computing 19(7), 2023.
\url{https://arxiv.org/abs/2111.02296}

\bibitem{buhrmantorenvliet1996}
H. Buhrman and L. Torenvliet,
\emph{P-Selective Self-Reducible Sets: A New Characterization of P},
Journal of Computer and System Sciences 53(2), 1996.
\url{https://doi.org/10.1006/jcss.1996.0062}
\end{thebibliography}
\end{document}
